% Section 2: Scope of Reproducibility
% Reproduce both Section 4 (Vision) and Section 5 (Language/Negation) of Pearce et al. (2024)

\section{Scope of Reproducibility}
\label{sec:scope}

We reproduce the main experimental sections of \citet{pearce2024bilinear}: Section 4 (Vision) demonstrating interpretable eigendecomposition on MNIST and Fashion-MNIST, and Section 5 (Language) discovering negation circuits in bilinear transformers using Sparse Autoencoder (SAE) analysis.

\subsection{Claims to Verify}

\subsubsection{Section 4: Vision (Image Classification)}

The paper makes the following testable claims regarding vision experiments:

\begin{enumerate}
    \item \textbf{Claim V1 (Low-Rank Emergence):} Regularization (input noise + weight decay) induces low-rank structure in the bilinear interaction tensor, as measured by effective rank.
    \item \textbf{Claim V2 (Interpretable Eigenvectors):} The eigenvectors of the regularized model's interaction tensor correspond to interpretable digit-like patterns.
    \item \textbf{Claim V3 (Overfitting Without Regularization):} Without regularization, the eigenspectrum is flat (high effective rank $\sim$150--200) and eigenvectors show noise-like overfitting patterns.
    \item \textbf{Claim V4 (Noise Importance):} Input noise augmentation is crucial for inducing interpretable structure, more so than weight decay alone.
    \item \textbf{Claim V5 (Cross-Seed Consistency):} Eigenvectors are consistent across random seeds (cosine similarity 0.8--0.9).
    \item \textbf{Claim V6 (Adversarial Generation):} Adversarial perturbations can be constructed directly from the eigendecomposition.
\end{enumerate}

\subsubsection{Section 5: Language (Negation Circuits)}

The paper discovers \textbf{sentiment negation circuits} in bilinear transformers:

\begin{enumerate}
    \item \textbf{Claim L1 (SAE Feature Discovery):} Sparse Autoencoders trained around MLP layers reveal interpretable features that can be analyzed via bilinear decomposition.
    \item \textbf{Claim L2 (Negation Features):} The model learns distinct features for ``not + negative words'' (e.g., ``not lost'') and ``not + positive words'' (e.g., ``not free'') that form opposing directions in feature space.
    \item \textbf{Claim L3 (Low-Rank Interactions):} 69\% of SAE output features achieve $>$0.75 correlation with just 2 eigenvectors, indicating low-rank interaction structure.
    \item \textbf{Claim L4 (Sparse Interaction Matrices):} The top 50 interactions between input and output SAE features form a small, interpretable submatrix.
    \item \textbf{Claim L5 (Cross-Model Generalization):} The same negation circuit structure appears in two different models (TinyStories and FineWeb-EDU).
\end{enumerate}

\subsection{Success Criteria}

\textbf{Vision (Section 4):}
\begin{itemize}
    \item Effective rank ratio (regularized / unregularized) $< 0.5$
    \item Top eigenvectors visually resemble digit/fashion templates
    \item Test accuracy within 2\% of paper values ($\sim$95\% with regularization)
    \item Cross-seed eigenvector cosine similarity $> 0.8$
\end{itemize}

\textbf{Language (Section 5):}
\begin{itemize}
    \item SAE features activate on expected negation patterns
    \item Negation feature pairs show opposing directions (negative cosine similarity)
    \item $>$60\% of features show $>$0.75 low-rank correlation
    \item Interaction submatrices are sparse and interpretable
\end{itemize}

\subsection{Paper Context}

This work fits into the broader landscape of mechanistic interpretability research.
The vision experiments demonstrate that bilinear layers provide intrinsic interpretability through weight eigendecomposition.
The language experiments show that this extends to understanding \emph{circuits}---computational subgraphs that implement specific behaviors like sentiment negation.
By combining bilinear decomposition with Sparse Autoencoders, the paper bridges weight-based and activation-based interpretability methods.
